%% LyX 2.4.4 created this file.  For more info, see https://www.lyx.org/.
%% Do not edit unless you really know what you are doing.
\documentclass[english]{article}
\usepackage[T1]{fontenc}
\usepackage[latin9]{inputenc}
\usepackage{enumitem}
\usepackage{amsmath}
\usepackage{amsthm}
\usepackage{geometry}
\geometry{verbose,tmargin=1in,bmargin=1in,lmargin=1in,rmargin=1in}

\makeatletter
%%%%%%%%%%%%%%%%%%%%%%%%%%%%%% Textclass specific LaTeX commands.
\numberwithin{equation}{section}
\numberwithin{figure}{section}
\newlength{\lyxlabelwidth}      % auxiliary length 

%%%%%%%%%%%%%%%%%%%%%%%%%%%%%% User specified LaTeX commands.
\usepackage{fancyhdr}
\usepackage{lastpage}
\pagestyle{fancy}
\usepackage{enumitem}
\usepackage{ifthen}
%sets math fonts to be more readable
\everymath=\expandafter{\the\everymath\displaystyle}
%\DeclareMathSizes{10}{10}{10}{10}
\usepackage{multicol}

\usepackage{tikz}
\usetikzlibrary{plotmarks}
\usepackage{pgfplots}
\pgfplotsset{compat=newest}

\usepackage{calc}
\usepackage[nomessages]{fp}

\usepackage{advdate}    % Advancing/saving dates
\usepackage{datetime}   % Dates formatting
\usepackage{datenumber} % Counters for dates
% Added by lyx2lyx
\setlength{\parskip}{\smallskipamount}
\setlength{\parindent}{0pt}

\makeatother

\usepackage{babel}
\begin{document}
\renewcommand{\author}{Dr.\,James Ashe}

\newcommand{\classNum}{137}

\newcommand{\classSec}{B}

\newcommand{\nameType}{Name:}

\newcommand{\examType}{Quiz 2.1-2.2}

\renewcommand{\date}{\formatdate{4}{3}{2013}}

%%First page's header%%%%%%%%%%%%%%%%%%%%%%
\fancypagestyle{firststyle} %%header settings for first pagr
{
	\fancyhead[L]{MTH \classNum{}(\classSec{}) \author{}\\
	\textbf{\examType{}} \date{}}
	\fancyhead[C]{\textbf{\nameType}}
	\setlength{\headsep}{14pt}
}
%%%%%%%%%%%%%%%%%%%%%%%%%%%%%%%%%%%%%%%%%%

%%Next page's header%%%%%%%%%%%%%%%%%%%%%%%
\setlist{nolistsep}
\lhead{\classNum{}(\classSec{})} 
\chead{\textbf{\nameType}} 
\cfoot{\thepage{}~of~\pageref{LastPage}} 
\setlength{\headsep}{5pt} 
%%%%%%%%%%%%%%%%%%%%%%%%%%%%%%%%%%%%%%%%%%%

%%Various settings; see the preamble for other settings%%%%%
%%\setenumerate[0]{leftmargin=12pt,itemindent=0pt,labelwidth=0pt}
\renewcommand{\labelenumi}{\textbf{\arabic{enumi}.}}
\renewcommand{\labelenumii}{\textbf{\alph{enumii}.}}
\thispagestyle{firststyle}
%%%%%%%%%%%%%%%%%%%%%%%%%%%%%%%%%%%%%%%%%%
\newcommand{\xmax}{10}
\newcommand{\xmin}{-10}
\newcommand{\xstep}{0} %must be xmin+n
\newcommand{\ymax}{10}
\newcommand{\ymin}{-10}
\newcommand{\ystep}{0} %must be ymin+n

\emph{Determine the domain and range of each relation, and state whether
the relation is a function. }
\begin{enumerate}[resume]
\item $f=\left\{ \left(10,1\right),\left(27,-1\right),\left(37,1\right),\left(10,2\right),\left(47,2\right)\right\} $\vfill\vspace{-2cm}
\item The curve shown in the graph below.\\
\renewcommand{\xmax}{8}
\renewcommand{\xmin}{-8}
\renewcommand{\xstep}{0} %must be xmin+n
\renewcommand{\ymax}{8}
\renewcommand{\ymin}{-8}
\renewcommand{\ystep}{0} %must be ymin+n
%%%%%%%
\begin{tikzpicture} 
\begin{axis}[height=4.5cm, 
	width=4.5cm, 
	axis lines=middle,
	minor x tick num={4},
	minor y tick num={4},
	%xtick={0,50,...,250},
	%ytick={0,10,20,...,40},
	grid=both,
	%axis line style={-latex}, 
	xmin=\xmin,xmax=\xmax, 
	ymin=\ymin,ymax=\ymax,
	%restrict y to domain=-1:10,
	%no markers, 
	xlabel={$x$},ylabel={$y$}, 
	%every axis x label/.append style={anchor=north}, 
	%every axis y label/.append style={anchor=east}
	]
	
	\addplot[color=red,solid,mark=*] coordinates {(-1,-3)};
	\addplot[mark=noner,red,thick,domain=\xmin:\xmax,samples=100,<->]{1/2(x+1)^2-3} node[anchor=135] {$$}; 
	
\end{axis}
\end{tikzpicture}\vspace{.5cm}
\item $x=y^{2}+2$. Choose the correct graph of the equation, determine
the domain and range, and determine whether $y$ is a function of
$x$.\\
%%%%%%%
\begin{tikzpicture} 
\begin{axis}[height=4.5cm, 
	width=4.5cm, 
	axis lines=middle,
	minor x tick num={4},
	minor y tick num={4},
	grid=both,
	%xtick={0,50,...,250},
	%ytick={0,10,20,...,40},
	%axis line style={-latex}, 
	xmin=\xmin,xmax=\xmax, 
	ymin=\ymin,ymax=\ymax,
	%restrict y to domain=-1:10,
	%no markers, 
	xlabel={$x$},ylabel={$y$}, 
	%every axis x label/.append style={anchor=north}, 
	%every axis y label/.append style={anchor=east}
	]
	
	\addplot[color=red,solid,mark=*] coordinates {(0,0)};
	\addplot[domain=\xmin:\xmax,thick,samples=100,draw=red,->]{sqrt(x)} node[right] {$$};
	\addplot[domain=\xmin:\xmax,thick,samples=100,draw=red,->]{-sqrt(x)} node[right] {$$};
	
\end{axis}
\end{tikzpicture}\hfill%%%%%%%
\begin{tikzpicture} 
\begin{axis}[height=4.5cm, 
	width=4.5cm, 
	axis lines=middle,
	minor x tick num={4},
	minor y tick num={4},
	%xtick={0,50,...,250},
	%ytick={0,10,20,...,40},
	grid=both,
	%axis line style={-latex}, 
	xmin=\xmin,xmax=\xmax, 
	ymin=\ymin,ymax=\ymax,
	%restrict y to domain=-1:10,
	%no markers, 
	xlabel={$x$},ylabel={$y$}, 
	%every axis x label/.append style={anchor=north}, 
	%every axis y label/.append style={anchor=east}
	]
	\addplot[color=red,solid,mark=*] coordinates {(2,0)};
	\addplot[domain=\xmin:\xmax,thick,samples=100,draw=red,->]{sqrt(x-2)} node[right] {$$};
	\addplot[domain=\xmin:\xmax,thick,samples=100,draw=red,->]{-sqrt(x-2)} node[right] {$$};
	
\end{axis}
\end{tikzpicture}\hfill%%%%%%%
\begin{tikzpicture} 
\begin{axis}[height=4.5cm, 
	width=4.5cm, 
	axis lines=middle,
	minor x tick num={4},
	minor y tick num={4},
	%xtick={0,50,...,250},
	%ytick={0,10,20,...,40},
	grid=both,
	%axis line style={-latex}, 
	xmin=\xmin,xmax=\xmax, 
	ymin=\ymin,ymax=\ymax,
	%restrict y to domain=-1:10,
	%no markers, 
	xlabel={$x$},ylabel={$y$}, 
	%every axis x label/.append style={anchor=north}, 
	%every axis y label/.append style={anchor=east}
	]
	
	\addplot[color=red,solid,mark=*] coordinates {(0,2)};
	\addplot[domain=\xmin:\xmax,thick,samples=100,draw=red,->]{sqrt(x)+2} node[right] {$$};
	\addplot[domain=\xmin:\xmax,thick,samples=100,draw=red,->]{-sqrt(x)+2} node[right] {$$};
	
\end{axis}
\end{tikzpicture}\hfill%%%%%%%
\begin{tikzpicture} 
\begin{axis}[height=4.5cm, 
	width=4.5cm, 
	axis lines=middle,
	minor x tick num={4},
	minor y tick num={4},
	%xtick={0,50,...,250},
	%ytick={0,10,20,...,40},
	grid=both,
	%axis line style={-latex}, 
	xmin=\xmin,xmax=\xmax, 
	ymin=\ymin,ymax=\ymax,
	%restrict y to domain=-1:10,
	%no markers, 
	xlabel={$x$},ylabel={$y$}, 
	%every axis x label/.append style={anchor=north}, 
	%every axis y label/.append style={anchor=east}
	]
	
	\addplot[color=red,solid,mark=*] coordinates {(0,2)};
	\addplot[domain=-2.45:2.45,thick,samples=100,draw=red,<->]{x^2+2} node[right] {$$};
	
\end{axis}
\end{tikzpicture}\vspace{2cm}
\end{enumerate}
\emph{Let $g(x)=-x^{2}+5x+3$.}
\begin{enumerate}[resume]
\item Find $g(-2)$\vspace{3cm}
\end{enumerate}
\emph{Use the difference quotient, $\frac{f(x+h)-f(x)}{h}$ for $h\neq0$,
to complete and simplify the following problem.}
\begin{enumerate}[resume]
\item $f(x)=x^{2}+1$.\vfill
\end{enumerate}

\end{document}
